% ===================================================================
% 研究の基本情報定義ファイル
%
% このファイルは研究計画書、説明文章、同意書、情報公開文書で共通して
% 使用される基本情報を定義します。
%
% 【重要】このファイルを編集するだけで、すべてのドキュメントに反映されます
% ===================================================================

% -------------------------------------------------------------------
% 研究タイトルと版管理
% -------------------------------------------------------------------
% ※ 以下はテンプレートの記載例として設定された架空（パロディ）の研究情報です。
%    実際に使用する際は、ご自身の研究の情報に置き換えてください。
\def\ResearchTitle{とてつもない予感の定量的評価に関する予備的観察研究}
\def\ProtocolVersion{第1.0版}
\def\ProtocolDate{202X年X月X日}

% -------------------------------------------------------------------
% 研究機関情報
% -------------------------------------------------------------------
\def\InstitutionName{京都大学医学部附属病院}
\def\InstitutionHead{京都大学医学部附属病院長}
\def\InstitutionFullName{京都大学大学院医学研究科・医学部及び医学部附属病院}
\def\DepartmentName{とてつもない予感情報学グループ}

% -------------------------------------------------------------------
% 研究代表者情報・研究事務局情報
% 【重要】これらの情報は private/members.yaml で管理されます
% -------------------------------------------------------------------
% 研究代表者、研究事務局、分担研究者等のメンバー情報は、
% private/members.yaml に記載し、scripts/compile_with_members.sh で
% コンパイルすることで自動的に各ドキュメントに転記されます。
%
% メンバー情報の設定方法：
% 1. cp private/members.yaml.example private/members.yaml
% 2. private/members.yaml を編集
% 3. ./scripts/compile_with_members.sh でコンパイル
%
% 詳細は README.md の「研究メンバー情報の管理とコンパイル」を参照してください。

% -------------------------------------------------------------------
% 研究期間
% -------------------------------------------------------------------
\def\StudyStart{研究機関の長の実施許可日}
\def\StudyEnd{202X年X月X日}
\def\EnrollmentStart{研究機関の長の実施許可日}
\def\EnrollmentEnd{202X年X月X日}
\def\DataAcquisitionStart{研究機関の長の実施許可日}
\def\DataAcquisitionEnd{202X年X月X日}

% -------------------------------------------------------------------
% 共同研究機関（必要に応じて使用）
% -------------------------------------------------------------------
\def\CollaboratingInstitution{共同研究機関名}
\def\CollaboratingPI{共同研究機関 研究責任者名}
